\documentclass[11pt,a4paper]{article}
\usepackage[margin=18mm]{geometry}
\usepackage[hidelinks]{hyperref}
\usepackage{parskip}
\pagestyle{empty}

\begin{document}

\noindent
Tze Hui Liew (corresponding author)\\
Faculty of Information Science and Technology, Multimedia University\\
Melaka, Malaysia\\
\href{mailto:thliew@mmu.edu.my}{thliew@mmu.edu.my}

\medskip
\noindent\today

\medskip
\noindent
The Editor-in-Chief\\
\emph{Journal of Systems and Software}

\medskip
\noindent
\textbf{Submission of a research paper:} ``It Is Not the Signature: An
Empirical Study of a Hidden Key-Parse Cost in JSON Web Token Validation''

Dear Editor,

We submit the above manuscript for consideration as a research paper in the
\emph{Journal of Systems and Software} (JSS).

The paper is an empirical study of performance attribution in a widely used
software library. Developers usually attribute the per-request cost of JSON
Web Token validation to cryptography and optimise or relocate it on that
basis. We show that in \texttt{jsonwebtoken}, a widely used Node.js library,
the signature is under a tenth of the validation cost. The main contributions
are:
\begin{itemize}
\item \textbf{Root cause.} Two independent decompositions locate 60\% of the
  cost in a discarded attempt to parse the shared secret as an asymmetric key,
  a failing parse used as a type test on every call.
\item \textbf{A transitive dependency sets its price.} Issued from C against
  eight identically built OpenSSL releases, the same calls cost
  394.4\,$\mu$s on 3.0.16 and 8.447\,$\mu$s on 3.5.8, so a base-image change
  alters the cost by more than an order of magnitude without any change to
  application code.
\item \textbf{Real-system and ecosystem evidence.} The OpenSSL~3.0 price
  persists in the Node.js that Ubuntu 24.04 ships; in a containerised service,
  cgroup CPU accounting attributes about half of the application's processor
  time per request to validation; among seven JWT libraries in four languages
  only this one uses a failing parse; and none of 327 sampled public call sites
  uses the pattern that avoids it.
\item \textbf{A safe fix.} The obvious optimisation fails the library's own
  key-confusion tests, because the parse is coupled to a security check. We
  give a narrow fix that passes the full test suite and makes validation 4.15
  times faster, and derive guidance for developers, library authors and
  performance engineers.
\end{itemize}

We believe the work suits JSS because it combines rigorous performance
measurement, root-cause analysis in a real library, a cross-ecosystem
comparison and a mining study of public code, and turns them into actionable
engineering guidance. In line with the journal's open science policy, every
number in the paper is generated from committed data in an open replication
package (doi:10.5281/zenodo.23011153), which also scripts the experiments that
need our testbed. The defect has been reported to the library maintainers
(auth0/node-jsonwebtoken issue \#1046); at the time of submission the proposed
fix had not been merged.

This manuscript is original, has not been published previously, and is not
under consideration for publication elsewhere. All authors have approved the
manuscript and agree with its submission. The authors declare no competing
interests. The use of generative AI in preparing the manuscript is declared in
the manuscript.

Thank you for considering our submission.

\medskip
\noindent
Yours sincerely,\\[4pt]
Tze Hui Liew, on behalf of all authors\\
Jannatul Ferdous Katha, Tasmia Tahmid Prova, Md Kishor Morol, Tze Hui Liew,
Dip Nandi

\end{document}
